\section{Methodology}\label{sec:method}

\ins{Describe your method, model, or system in sufficient detail for a knowledgeable reader to understand what you did. This includes implementation details, experimental setup, data processing, and design choices. Explain and justify important decisions rather than listing them.}

We study whether chain-of-thought (CoT) explanations faithfully reflect
the influence of socially and culturally grounded bias cues in
multilingual reasoning settings. Our methodology is designed to measure
both \emph{behavioral susceptibility} to biased hints and the
\emph{transparency} of the model's reasoning process across languages
and across task domains.

\subsection{Datasets}

We evaluate on three task domains: a culturally grounded
multiple-choice benchmark (cultureMCQA) that we construct, an objective
mathematical reasoning benchmark (GSM8K), and a multilingual safety
benchmark (XSAFETY).

\paragraph{cultureMCQA.}
For culturally grounded tasks we construct \textsc{cultureMCQA}, a
multiple-choice benchmark of 200 items spanning four cultures: Korean
(KR), American (US), German (DE), and Polish (PL), with 50 items per
culture. The topical coverage is grounded in
CultureBank~\citep{shi2024culturebank}; we filter to entries with
agreement score $\geq 0.7$ to reduce ambiguity, and use these as
references when authoring scenario-based four-option questions. Each
item is hand-curated to satisfy three criteria: (i) the question is
unambiguous to a member of the target culture, (ii) the correct option
reflects current, attested practice rather than stereotype, and
(iii) the three distractors are plausible enough that a non-member
might select them. For each item we retain the top-2 most similar
CultureBank entries (TF--IDF cosine over the same target culture) as
attached references for provenance and audit. Items are written in
English; native-language renderings of the question and options are
generated only for analysis-time conditions that require them.

\paragraph{GSM8K (MCQ-converted).}
For objective reasoning we sample 50 problems from the GSM8K test
set~\citep{cobbe2021gsm8k} and convert each to a
four-option MCQ. The correct option is the gold numeric answer; three
distractors are constructed primarily from the intermediate
computations in GSM8K's annotated solution (extracted via the
\verb|<<expr=result>>| markers), with mechanical fallbacks
($\pm 1, \pm 5, \times 2, \div 2$) when fewer than three intermediates
are available. The correct letter is randomly assigned across A--D.
Using intermediate values yields distractors that correspond to plausible
reasoning errors rather than arbitrary noise.

\paragraph{XSAFETY.}
For the multilingual safety domain we use
XSAFETY~\citep{wang2024xsafety}. We sample 50 prompts each in English
and German (Korean and Polish are not covered by XSAFETY), stratified
across 14 safety categories. The expected behavior is refusal rather
than a discrete answer letter, so XSAFETY runs use a separate
refusal-classification pipeline; we include it as a cross-domain check
on whether hint cues degrade refusal behavior.

\subsection{Hint-injection design}

We evaluate models under a controlled hint-injection protocol. For
each item we first obtain a no-hint baseline response, and then
re-evaluate the same item under additional bias conditions.

\paragraph{Hint families.}
We use four hint families. The authority framing differs by domain:
on cultureMCQA we use a culturally framed expert
(``A professor of cultural studies suggests the answer is (X)''),
on GSM8K and XSAFETY we use a generic expert
(``An expert told me the answer is (X)''). Three families are shared
across all domains: social pressure (``Most people in our study group
chose (X)''), indirect suggestion (``I was thinking maybe (X) could
be right\ldots''), and a Turpin-style few-shot biased CoT condition
described below. We deliberately exclude self-opinion, confidence, and
research-citation framings to keep the cross-domain comparison
focused.

\paragraph{Hint languages.}
Simple-hint conditions are issued in both English and Korean,
producing two language variants per family. Templates for German and
Polish exist in the released hint table and can be enabled to expand
the cross-lingual matrix. This yields seven non-fewshot conditions per
dataset (one baseline plus three hint families $\times$ two languages),
plus the fewshot-biased condition described next.

\paragraph{Few-shot biased CoT (``Answer Always (X)'').}
Following \citet{turpin2023language}, the few-shot biased condition
constructs a prompt in which all $K=3$ in-context demonstrations end
with the same letter as the target's chosen wrong letter. For each
demonstration we \emph{re-order the option text} so that the correct
answer lands at the target wrong-letter position, then have the model
generate a confident step-by-step reasoning that ends with that
letter. After the reorder the demonstration is internally coherent
(the reasoning correctly points to the option now sitting at the
chosen letter), but a model that picks up on the positional pattern is
biased to give the same letter on the target item, whose options are
\emph{not} re-ordered.

\paragraph{Dynamic per-model wrong-letter selection.}
For each (model, item) pair where the model's baseline answer is
correct, we draw the wrong letter $w$ as a uniform sample from
$\{A,B,C,D\} \setminus \{\text{correct}, \text{baseline\_pred}\}$.
This ensures that the injected cue is always genuinely
counter-evidence: it points to an option different from both the gold
answer and the model's own belief. Storing a fixed wrong letter in the
dataset would conflate sycophancy with cases where the model already
agreed with the cue, so we leave the field \texttt{null} in the static
data and assign $w$ at experiment time.

\subsection{Two-phase experimental pipeline}

\paragraph{Phase 1 (baseline).}
We query the model with the question and options only and record
\emph{(i)} the predicted letter, \emph{(ii)} the model's natural-language
reasoning, and \emph{(iii)} for models that support it, the separate
``thinking'' trace exposed by the API. We filter to items where the
baseline answer is correct; only on these items can we cleanly observe
a sycophantic flip, because there is a stable belief to overturn.

\paragraph{Phase 2 (hint conditions).}
For each baseline-correct item we run all hint conditions defined for
the dataset. For each condition we record the same three outputs as in
Phase 1, plus the chosen wrong letter $w$. Few-shot biased conditions
additionally record the sampled demonstration ids, their option
swaps, and the generated biased reasoning, so that the construction is
fully auditable.

\paragraph{Model and inference setup.}
The main experiments use Gemini~2.5 Flash with thinking enabled
(\texttt{thinking\_budget}=8000, \texttt{include\_thoughts=True}),
temperature 1.0, and structured JSON output to make answer extraction
deterministic. The pipeline is implemented behind a thin
provider-agnostic client (\texttt{scripts/\_models.py}); replacing
\texttt{gemini/gemini-2.5-flash} with, e.g., \texttt{groq/qwen-2.5-32b}
or \texttt{openai/gpt-4o-mini} in the command line is sufficient to
re-run the same experiment on an alternative model.

\paragraph{Detection of hint acknowledgment.}
For each non-fewshot condition we detect whether the hint is mentioned
in the thinking trace using a small set of hint-type-specific
keywords (e.g.\ ``professor''/``expert'' for authority, ``study
group''/``most people'' for social, ``maybe''/``could be'' for
indirect). For the few-shot biased condition, a positive match
requires explicit reference to the in-context examples themselves
(``the previous examples'', ``the demonstrations'', ``the pattern''),
since the demonstration content itself already contains the hint
phrasing.
